\documentclass[11pt,a4paper,fontset=none]{ctexart}
\usepackage[margin=2.2cm]{geometry}
\usepackage{newtxmath}
\setmainfont{Times New Roman}
\setsansfont{Times New Roman}
\newcommand{\KaiPath}{/System/Library/AssetsV2/com_apple_MobileAsset_Font8/88d6cc32a907955efa1d014207889413890573be.asset/AssetData/}
\setCJKmainfont{Kaiti.ttc}[Path=\KaiPath, FontIndex=0, BoldFont=Kaiti.ttc, BoldFeatures={FontIndex=3}]
\setCJKsansfont{Kaiti.ttc}[Path=\KaiPath, FontIndex=0, BoldFont=Kaiti.ttc, BoldFeatures={FontIndex=3}]
\setCJKmonofont{Kaiti.ttc}[Path=\KaiPath, FontIndex=0]
\usepackage{amsmath}
\usepackage{booktabs,tabularx,array,multirow,longtable}
\usepackage{enumitem}
\usepackage{xurl}
\usepackage[hidelinks]{hyperref}
\usepackage{fancyhdr}
\setlist{nosep,leftmargin=1.4em}
\ctexset{section={format=\Large\bfseries\raggedright},subsection={format=\large\bfseries\raggedright}}
\setlength{\parskip}{0.3em}\setlength{\emergencystretch}{2em}
\renewcommand{\arraystretch}{1.2}
\newcolumntype{L}{>{\raggedright\arraybackslash}X}
\pagestyle{fancy}\fancyhf{}
\lhead{\small BAI AI Lab 代码能力评测}\rhead{\small 2026-10-27}\cfoot{\small\thepage}
\renewcommand{\headrulewidth}{0.4pt}

\begin{document}

\begin{center}
{\LARGE\bfseries 阶梯代码评测集设计说明}\\[0.4em]
{\large L1--L5 五档难度：正确性、效率与代码质量}\\[0.3em]
{\small 2026-10-27 \quad codingBench \texttt{cb ladder}}
\end{center}

\section*{摘要}
现有的 LiveCodeBench 最新 35 题已被一线模型做满：复核判分后 DeepSeek-V4.1-Flash 与 DeepSeek-V4-Pro 均为 35/35，Kimi K3.1 Pro 为 34/35，正确率已无法区分模型；同一批答案的运行时间却相差明显（相对最快答案的几何平均倍数为 1.29、2.00、2.74）。为此我们构建了一套按人类难度分为五档的测试集：
\begin{itemize}
  \item \textbf{L1--L4}：USACO 2024--25 与 2025--26 两个赛季的铜、银、金、铂金组，L3、L4 另加入 2026 年的 EGOI 与 IOI 新题；
  \item \textbf{L5}：FrontierCS 算法赛道，没有已知最优解，按 0--100 给连续分。
\end{itemize}
所有题目、测试数据和题解都来自官方发布。每道题同时报告四类结果：
\begin{itemize}
  \item 宽时限下的正确性；
  \item 官方时限下的正确性；
  \item 相对官方解的运行时间与内存；
  \item 相对官方解的时间、空间复杂度增长阶（是否用了更优或更差的解法）；
  \item Python 档位的静态代码质量。
\end{itemize}
评测流程已实现为一条命令（\texttt{cb ladder}）：接入模型 API 后自动生成、判题、复测性能并输出报告。

\section{背景与目标}
\begin{enumerate}
  \item \textbf{恢复区分度。}题目要足够难，并且难度分层：弱模型在低档拉开差距，强模型在高档拉开差距。
  \item \textbf{不止看对错。}同时测“够不够快”，即官方时限下能否通过、运行时间是官方解的多少倍，以及内存、代码质量、调用成本与时延。
  \item \textbf{防止训练数据泄露。}记录每道题的比赛日期，并单独报告 2026 年新题的成绩作为对照。
  \item \textbf{可复现。}只用官方测试数据和官方题解，来源钉在固定的版本上，判题规则与官方一致。
\end{enumerate}

\section{设计原则}
\begin{itemize}
  \item \textbf{难度用人类标尺划分。}USACO 的四个组别由真实选手的晋级线划定（例如 2026 年第一场金组得分 800 以上晋级铂金），EGOI、IOI 是国际奥赛，所以档位对应的是人类选手的水平，而不是哪个模型觉得难。
  \item \textbf{只用官方数据。}AtCoder 自 2024 年 12 月起停止公开测试数据，其 robots.txt 也禁止抓取提交记录，难题的官方题解里多半没有代码；LeetCode 的测试不公开，用户协议禁止批量抓取。这两个来源都给不出可信的标准答案，因此没有采用。
  \item \textbf{高档用连续分。}USACO 按测试点给部分分，EGOI、IOI 按子任务给分，FrontierCS 按 checker 给 0--100 的连续分，高档位不会像通过率那样很快饱和。
  \item \textbf{正确与速度分开测。}每个测试点只运行一次，时限放宽到官方的 5 倍，再按实际 CPU 时间同时得出两种判定：宽时限判的是算法对不对，官方时限判的是够不够快。
  \item \textbf{先验证判题本身。}建库时先用我们的判题程序跑官方题解；没有任何一份官方解能全部通过的题直接剔除，这同时检验了测试数据、比对规则和判题程序。
\end{itemize}

\section{题库组成}
\begin{center}
\small
\begin{tabularx}{\linewidth}{@{}l>{\raggedright}p{4.2cm}rlll@{}}
\toprule
档位 & 来源 & 题数 & 比赛时间 & 作答语言 & 计分 \\
\midrule
L1 基础 & USACO 铜组 & 21 & 2024-12 至 2026-02 & Python & 测试点等分 \\
L2 中等 & USACO 银组 & 21 & 2024-12 至 2026-02 & Python & 测试点等分 \\
\multirow{2}{*}{L3 困难} & USACO 金组 & 21 & 2024-12 至 2026-02 & C++ & 测试点等分 \\
 & EGOI 2026（2026 子集） & 8 & 2026-05 & C++ & 子任务（取最低） \\
\multirow{2}{*}{L4 顶尖} & USACO 铂金组 & 21 & 2024-12 至 2026-02 & C++ & 测试点等分 \\
 & IOI 2026 批处理题（2026 子集） & 2 & 2026-08 & C++ & 子任务（取最低） \\
L5 开放式 & FrontierCS 算法赛道（按题号均匀抽样） & 30 & 题库 v1.1 & C++17 & 0--100 连续分 \\
\bottomrule
\end{tabularx}
\end{center}
\begin{itemize}
  \item USACO 每个赛季有若干场比赛，每场每组 3 题。2024--25 赛季取 12 月、1 月、2 月和美国公开赛共 4 场，2025--26 赛季取 1 月、1 月底和 2 月 3 场线上赛（3 月 28 日的公开赛为线下监考赛，官网上没有它的题目和测试数据），合计 84 题，四个组别各 21 题。新赛季开始后可以一键补入。
  \item IOI 2026 有 6 题，其中 4 题是“通信”题，需要管理程序与多个选手进程交互，现在的判题程序还不支持，暂未列入；先纳入 2 道批处理题（Monuments、Magic City）。
  \item FrontierCS 题库有 180 道开放式难题，固定抽 30 题。题库没有提供类别标签，所以按题号均匀抽取（第 0、6、13、23……310 号），需要时可以跑全部 180 题。
  \item 表中题数是建库前的计划数。建库校验时，没有官方解能通过的题，以及有多个合法答案却没有 checker 的题，都会被剔除，最终以 \texttt{samples/ladder.json} 为准。
  \item 原始数据共约 4.7 GB。L1--L4 的题面、测试数据与官方题解约 2.6 GB，FrontierCS 约 2.2 GB。题面和测试数据只存在服务器上，不进仓库、不公开；仓库里只记录题号、档位、日期和时限。
\end{itemize}

\section{判题与计分原理}
\subsection{运行环境}
判题在服务器本机进行，不依赖 Docker。
\begin{itemize}
  \item \textbf{运行身份与限制。}答案以无特权用户 \texttt{bench} 运行，经一个小的启动器设置五项上限：
    \begin{itemize}
      \item 地址空间：内存上限加 64 MB；
      \item 栈：等于内存上限；
      \item CPU 时间；
      \item 输出大小：256 MB；
      \item 进程绑定在一个独占的核上。
    \end{itemize}
  \item \textbf{计时。}CPU 时间和峰值内存取自子进程本身的资源统计（\texttt{wait4}）。
  \item \textbf{Python 档。}用 CPython 3.11，只允许标准库。官方时限是 USACO 给 Python 的时限，默认 4 秒，是 C++ 的两倍；题目注明加时的，按比例放大。
  \item \textbf{C++ 档。}编译命令为 \texttt{g++ -O2 -std=gnu++17}。IOI 题与官方 grader 一起编译，参数与 IOI 2026 的 CMS 一致（\texttt{-DEVAL -std=gnu++20 -O2}，不静态链接）。
\end{itemize}

\subsection{判定与比对}
\begin{itemize}
  \item 每个测试点的结果是以下之一：
    \begin{itemize}
      \item 正确（含部分分）；
      \item 答案错误；
      \item 超时；
      \item 超内存；
      \item 运行错误；
      \item 编译错误。
    \end{itemize}
    宽时限为官方时限的 5 倍（\texttt{LADDER\_LENIENT}）。在宽时限内通过、但 CPU 时间超过官方时限的测试点，在“官方时限”口径下记 0 分。
  \item \textbf{USACO：}忽略空白、逐个 token 比对，与 USACO 的规则一致。
  \item \textbf{EGOI：}用官方 checker（CMS 约定，checker 输出 0--1 的得分）。
  \item \textbf{IOI：}官方 grader 与选手代码一起编译，输出交给官方 checker。
  \item \textbf{FrontierCS：}用官方 testlib checker 或交互器。规则照搬官方的 \texttt{judge\_engine.js}：
    \begin{itemize}
      \item 只解析 checker 自己输出的 \texttt{Ratio}；
      \item 退出码 0 表示正确，7 表示部分分；
      \item 达到参考解得 60 分、超过参考解按对数加分，这部分写在各题的 checker 里。
    \end{itemize}
\end{itemize}

\subsection{计分}
\begin{itemize}
  \item \textbf{每题得分（0--100）}按来源的官方规则计算（公式见第 5 节）：
    \begin{itemize}
      \item USACO：每个测试点分值相同；
      \item EGOI、IOI：子任务取其中所有测试点的最低分，再乘以该子任务的分值（CMS GroupMin）；
      \item FrontierCS：所有测试点 Ratio 的平均值乘以 100。
    \end{itemize}
  \item \textbf{解出}指每个测试点都满分。
  \item \textbf{采样}：每档的采样次数为 L1、L2、L5 各 1 次，L3、L4 各 3 次（\texttt{LADDER\_N}），temperature 0.2。高档题量少，多采样能减小单次运气的影响。pass@1 是所有采样中解出的比例；pass@$n$ 是至少有一次解出的题所占比例。
  \item \textbf{提示词}：USACO、EGOI 沿用 LiveCodeBench 的写法，并给出时限和内存限制；IOI 要求只实现规定的函数；FrontierCS 用官方生成脚本的原始提示词，便于和它的公开排行榜对照。
\end{itemize}

\subsection{性能对照}
建库时，每题的所有官方题解都要在全部测试点上判一遍；能全部通过的官方解中，最快的一份作为该题的性能对照。

评测时，对每份解出的答案：
\begin{itemize}
  \item 取官方解最慢的 3 个测试点；
  \item 答案与官方解交替各运行 3 次，取 CPU 时间的中位数；
  \item 报告“答案 CPU 时间 ÷ 官方解 CPU 时间”，各档取几何平均。
\end{itemize}
只有官方解在测试点上至少需要 0.05 秒 CPU 的题才做这项对照，否则时间主要是进程启动开销。Python 档只与 Python 官方解比较，没有 Python 官方解的题不做这项对照。

\subsection{时间与空间复杂度}
正确率只说明“做没做对”，复杂度还能说明“用了什么样的解法”。一份答案如果增长得比官方解还慢，往往意味着更好或更新颖的算法；增长得更快，则说明它虽然在这批测试上通过了，但算法在渐近意义上更差。

官方测试并不是按规模排列的序列：在 USACO 里，耗时达到可测量程度的测试点几乎都是最大规模，输入大小相差不到 2 倍，直接按输入大小拟合基本不可行（首轮评测中 109 份解出的答案只有 20 份可测）。因此以官方解自己的工作量为横轴：
\begin{itemize}
  \item 选官方解至少需要 2 ms CPU（扣除空程序）的测试点，最多 20 个，均匀覆盖其耗时范围；
  \item 答案与官方解在这些测试点上交替各运行 5 次，取 CPU 时间和峰值内存的中位数，压低毫秒级的测量噪声；
  \item 拟合 $\log y_{\text{答案}} = a + k \log y_{\text{官方}}$。若官方解为 $\Theta(f(n))$、答案为 $\Theta(f(n)^k)$，$k$ 就是这个指数；
  \item 只在拟合优度 $R^2 \ge 0.8$ 时下结论：$k \le 0.75$ 为更优，$0.75 < k < 1.33$ 为同阶，$k \ge 1.33$ 为更差；否则记为不确定。
\end{itemize}
这些测试点的输入大小跨度达到 8 倍时，另外报告两者按输入大小的增长幂次（参照 BigO(Bench)\cite{bigo} 的规模扫描）。增长慢于官方解的答案单独列出，供人工复核。下一步为 EGOI 用官方生成器、为 USACO 编写生成器，造出不同规模的输入，做真正的规模扫描。

% generated from tools/ladder_report.py CRITERIA: the same table as in every cb ladder report
\section{评测维度与打分标准}
各维度分开报告，不合成总分；表格中 ↑ 越大越好，↓ 越小越好，没有箭头的只作描述。
{\small\begin{longtable}{@{}p{1.6cm}p{2.6cm}p{10.3cm}p{0.9cm}@{}}
\toprule 维度 & 指标 & 计算方法 & 方向 \\ \midrule \endhead
正确性 & 测试点得分 $s_t$ & 每个测试点 0–1：USACO 忽略空白逐 token 比对，相同为 1、否则 0；EGOI、IOI 取官方 checker 给出的 0–1 分；FrontierCS 取官方 checker 输出的 Ratio（达到参考解为 0.6，超过参考解按对数加分，最高 1）。程序编译失败、运行出错、超时或超内存，该测试点为 0。 &  \\
正确性 & 题目得分 $S$（0–100） & USACO：$S = 100 \cdot \frac{1}{T}\sum_t s_t$（每个测试点等分）；EGOI、IOI：$S = 100 \cdot \frac{\sum_g w_g \min_{t \in g} s_t}{\sum_g w_g}$（子任务 $g$ 取其中测试点的最低分，再乘以子任务分值 $w_g$）；FrontierCS：$S = 100 \cdot \frac{1}{T}\sum_t \mathrm{Ratio}_t$。 & ↑ \\
正确性 & 解出 & 所有测试点 $s_t = 1$。 &  \\
正确性 & pass@1、平均得分、pass@$n$ & pass@1 = 解出的采样数 ÷ 采样总数 × 100；平均得分 = 所有采样的 $S$ 的平均；pass@$n$ = 至少一次采样解出的题数 ÷ 题数 × 100。每档采样次数：L1、L2、L5 各 1 次，L3、L4 各 3 次，temperature 0.2。 & ↑ \\
正确性 / 效率 & 宽时限与官方时限 & 每个测试点只运行一次，时限为官方时限的 5 倍（宽时限）。宽时限口径：按上面的规则计分。官方时限口径：CPU 时间超过官方时限的测试点记 0，再按同样的规则算 $S$、pass@1。前者回答“算法对不对”，后者回答“是否和人类选手一样在时限内通过”。 & ↑ \\
效率 & 运行时间倍数 & 只对解出的答案、且官方解在某个测试点上至少需要 0.05 秒 CPU 的题：取官方解最慢的 3 个测试点，答案与官方解交替各运行 3 次取中位数，$r$ = 答案 CPU 时间 ÷ 官方解 CPU 时间（取官方解最慢的那个测试点）；各档报告 $r$ 的几何平均和题数。同语言比较。 & ↓ \\
效率 & 内存倍数 & 解出的答案在所有测试点上的峰值内存 ÷ 官方解的峰值内存，各档取几何平均。 & ↓ \\
复杂度 & 增长幂次 $k$ 与 $\Delta k$ & 以测试点输入字节数为规模 $x$，CPU 时间（或峰值内存）扣除空程序开销为 $y$，在 $y$ 高于噪声下限（10 ms / 2 MB）的测试点上拟合 $\log y = a + k\log x$；在答案与官方解都可测的测试点上拟合 $\log(y_{\text{答案}}/y_{\text{官方}})$ 对 $\log x$ 的斜率 $\Delta k$。$\Delta k \le -0.3$ 为更优（增长慢于官方解，列出供人工复核），$|\Delta k| < 0.3$ 为相同，$\Delta k \ge 0.3$ 为更差；输入大小变化不足 8 倍或可测测试点少于 4 个记为不可测。时间、空间分别判定。 &  \\
失败类型 & 每份未解出答案的主要原因 & 按顺序取第一个成立的：无代码、编译错误、答案错误、运行错误、超内存、超时、部分分（部分测试点得分）；宽时限下解出但官方时限下没有全部通过的记为“只在官方时限下超时”。 &  \\
静态质量 & Pylint、Radon、SonarQube（只限 L1–L2 的 Python 答案） & 按 arXiv:2610.19167 的口径与版本：Pylint 3.3.5 默认配置的问题数（C/R/W/E）与 0–10 分，Radon 6.0.1 的圈复杂度（所有条目之和）、可维护性指数 MI、SLOC，SonarQube 10.4 的可维护性/可靠性/安全问题数；报告正确答案的每份平均与每百行 SLOC 的问题数。 & 见表头 \\
成本与时延 & token、费用、延迟、首个 token 时间 & 费用 = (输入 − 缓存命中) × 输入单价 + 缓存命中 × 命中单价 + 输出（含思考）× 输出单价，单价按每百万 token，分时计价按每次调用的开始时间；延迟为每次调用从发出到收完的时间。 & ↓ \\
泄露检查 & 2026 子集与 USACO 子集 & L3、L4 中 EGOI/IOI 2026 新题与同档 USACO 旧题分开报告 pass@1 与平均得分；新题明显更差时提示旧题可能见过。只作提示，不作结论。 &  \\
\bottomrule
\end{longtable}}


\section{预期结果}
以下是根据已有数据作出的预测，不是测量结果；它们的作用是给出判断标尺，实际数字以评测为准。
\begin{center}
\small
\begin{tabularx}{\linewidth}{@{}l>{\raggedright}p{3.6cm}L@{}}
\toprule
档位 & 一线模型的预期 & 依据 \\
\midrule
L1 & pass@1 95\% 以上 & 难度与 LiveCodeBench 中等题相当，现有模型在 LCB 上已做满 \\
L2 & 80--100\%；官方时限下略低 & Python 运行较慢，USACO 提示高组别的题 Python 未必能在时限内通过 \\
L3 & USACO 金组 50--85\%；EGOI 更低 & EGOI 的最后几题接近 IOI 难度 \\
L4 & pass@1 15--50\%，平均得分 40--75 & 2026 年第一场铂金组只有 2 名认证选手超过 500/1000 分；OpenAI 的推理系统在 IOI 2025 线上赛达到金牌分数、与人类同场排名第 6 \cite{oaiioi}，但那是 50 次提交、有评测反馈的条件，我们是单次作答、没有反馈 \\
L5 & 40--57 分 & FrontierCS 2026-10-18 公布的 6 个一线模型为 47.8--56.5 分，人类专家 86.99 \cite{fcs} \\
性能 & 运行时间约为官方解的 1--3 倍 & 此前 LCB 35 题中三个模型相对最快答案为 1.29--2.74 倍 \\
\bottomrule
\end{tabularx}
\end{center}
\begin{itemize}
  \item \textbf{区分度}主要来自 L3--L5：低档用来确认基本能力，同时测小模型。
  \item \textbf{两种时限的差距}预计在 Python 档（L2）最明显，C++ 档较小。如果某模型在宽时限下解出、官方时限下失败的题特别多，说明它常写出思路正确但复杂度偏高的解法。
  \item \textbf{如果 USACO 子集明显好于同档的 2026 子集}，提示旧题可能出现在训练数据里。但两部分的难度分布不完全相同，只能作为提示，不能作为结论。
  \item \textbf{成本}：一次完整评测约 228 次调用，包括 L1 21 次、L2 21 次、L3 87 次、L4 69 次、L5 30 次。按此前 LCB 难题的单次费用粗估：
    \begin{itemize}
      \item DeepSeek-V4.1-Flash 约 \$2.5；
      \item DeepSeek-V4-Pro 约 \$22；
      \item Kimi K3.1 Pro 约 \$56。
    \end{itemize}
    高档题思考更长，实际可能更高。
\end{itemize}

\section{局限与后续工作}
\begin{itemize}
  \item \textbf{题目时效。}USACO 题目截至 2026 年 2 月，可能已进入部分模型的训练数据，所以记录了比赛日期并设 2026 子集作对照。USACO 2026--27 赛季公布后补入新题。
  \item \textbf{IOI 覆盖不全。}IOI 2026 的 4 道通信题暂不支持，下一步实现“管理程序 + 多个选手进程”的判题。
  \item \textbf{C++ 档没有静态质量表。}所用工具只支持 Python，SonarQube 10.4 社区版不分析 C/C++。
  \item \textbf{时限照搬官方。}评测机与 USACO、EGOI、IOI 官方评测机的速度不同。时限沿用官方值，建库时会确认官方解在这台机器上能在官方时限内通过。
  \item \textbf{L5 判题程序待验证。}FrontierCS 官方用 go-judge 容器判题，我们按其规则实现了本机判题程序。下一步重判它仓库里公开的历史解，与公开排行榜核对。
  \item \textbf{AlgoTune 性能专项未上线。}它测“把数值程序写快”，依赖约 3--4 GB 的软件包，等服务器网络恢复后再上线。
  \item \textbf{服务器网络很慢。}服务器入口带宽实测约 350 KB/s，第一跳网关就有 25\% 丢包，首次下载全部数据约需 4 小时。
\end{itemize}

\section{使用方法}
\begin{verbatim}
cb ladder-data                # 后台下载（可断点续传）；--status 看进度
cb ladder-build               # 整理题目；--calibrate 判官方解并选出性能对照
export MODEL=... OPENAI_BASE_URL=... OPENAI_API_KEY=...
cb ladder --smoke             # 每档 1 题，先确认接口
cb ladder                     # 全量：生成 -> 判题 -> 性能复测 -> 复杂度 -> 静态质量 -> 报告
cb ladder-report <运行A> <运行B> --out <目录>   # 多模型对比报告
\end{verbatim}
管路的测试 \texttt{tests/test\_ladder.py} 用合成题和一个假的模型接口，把生成、判题（四种判题方式、超时、超内存、两种时限）、性能复测和报告完整跑一遍，在 CI 和服务器上均已通过。

\begin{thebibliography}{9}\small
\bibitem{usaco} USACO 2024--25、2025--26 赛季比赛结果、题面、测试数据与题解. \url{https://usaco.org/}
\bibitem{egoi} EGOI 2026 题目包. \url{https://github.com/egoi-org/egoi-2026}（commit 08eedd3）
\bibitem{ioi} IOI 2026 题目档案（CC BY 4.0）. \url{https://github.com/ioi/task-archive}（commit 70278df）
\bibitem{fcs} Q. Mang, W. Chai, Z. Li, et al. FrontierCS: Evolving Challenges for Evolving Intelligence. arXiv:2512.15699, 2025. 数据与结果：\url{https://github.com/FrontierCS/Frontier-CS}（commit 12ab268）
\bibitem{lcb} N. Jain, K. Han, A. Gu, et al. LiveCodeBench: Holistic and Contamination Free Evaluation of Large Language Models for Code. arXiv:2403.07974, 2024.
\bibitem{usacobench} Q. Shi, M. Tang, K. Narasimhan, et al. Can Language Models Solve Olympiad Programming? arXiv:2404.10952, 2024.
\bibitem{enamel} R. Qiu, W. W. Zeng, J. Ezick, et al. How Efficient is LLM-Generated Code? A Rigorous \& High-Standard Benchmark. arXiv:2406.06647, 2024.
\bibitem{perfbench} N. M. Le, Y. Xu, Z. Wang, et al. Rethinking Code Performance Benchmarks for LLMs. arXiv:2607.07619, 2026.
\bibitem{quality} J. Huang, T. Mahmud, M. Che, et al. Beyond Functional Correctness: Static Code Quality Profiling of LLM-Generated Code. arXiv:2610.19167, 2026.
\bibitem{bigo} P. Chambon, B. Rozière, B. Sagot, G. Synnaeve. BigO(Bench): Can LLMs Generate Code with Controlled Time and Space Complexity? arXiv:2503.15242, 2025.
\bibitem{algotune} O. Press, B. Amos, H. Zhao, et al. AlgoTune: Can Language Models Speed Up General-Purpose Numerical Programs? arXiv:2507.15887, 2025.
\bibitem{oaiioi} OpenAI. 关于 IOI 2025 线上赛成绩的公告, 2025-08-11. \url{https://x.com/OpenAI/status/1954969035713687975}
\end{thebibliography}

\end{document}
